% Chapter 6: Technologies, Stacks and Tools Used
\chapter{Technologies, Stacks and Tools Used}
This chapter provides a comprehensive overview of all technologies, frameworks, libraries, hardware components, and development tools employed in the Intelligent Laundry System project. Each technology is briefly described along with its specific role in the system.


\section{Development Environment}
The entire project was developed on a Windows 11 machine, with cross-platform compatibility ensured for the mobile application.

\begin{table}[H]
    \centering
    \begin{tabular}{|l|l|}
        \hline
        \textbf{Component} & \textbf{Specification} \\
        \hline
        Operating System & Windows 11 (Development), Android/iOS (Deployment) \\
        \hline
        Code Editors & Visual Studio Code, Arduino IDE \\
        \hline
        Version Control & Git, GitHub \\
        \hline
        Package Managers & pip (Python), npm (Node.js) \\
        \hline
    \end{tabular}
    \caption{Development Environment}
    \label{tab:dev_env}
\end{table}

\section{Programming Languages}
\begin{table}[H]
    \centering
    \begin{tabular}{|l|c|l|}
        \hline
        \textbf{Language} & \textbf{Version} & \textbf{Usage} \\
        \hline
        Python & 3.9+ & Backend API development (FastAPI), AI integration, data processing \\
        \hline
        C++ & Arduino compatible & Arduino firmware for hardware control \\
        \hline
        TypeScript & 4.9+ & React Native mobile application (type-safe development) \\
        \hline
        JavaScript & ES6 & React Native components and logic \\
        \hline
    \end{tabular}
    \caption{Programming Languages}
    \label{tab:langs}
\end{table}

\section{Backend Technologies}

\subsection{FastAPI}
FastAPI is a modern, high-performance web framework for building APIs with Python. It was chosen for its:
\begin{itemize}
    \item Automatic interactive API documentation (Swagger UI)
    \item Fast execution speed (on par with Node.js and Go)
    \item Built-in data validation using Pydantic
    \item Asynchronous support for handling multiple requests
\end{itemize}
Two independent FastAPI applications were developed:
\begin{itemize}
    \item \textbf{API 1 (Port 8000):} Clothing Material Identifier – accepts images and returns fabric analysis
    \item \textbf{API 2 (Port 8001):} Washing Parameter Predictor – accepts fabric data and returns wash parameters
\end{itemize}

\subsection{Uvicorn}
Uvicorn is an ASGI server used to run the FastAPI applications during development and testing. It provides high performance and supports automatic reloading for rapid iteration.

\subsection{Pydantic}
Pydantic is a data validation library that enforces type hints at runtime. It was used to define request and response models, ensuring that data exchanged between APIs and the mobile app strictly adheres to expected formats (e.g., \texttt{dirt\_level} is an integer between 1 and 5).

\subsection{Python-dotenv}
This library loads environment variables from a \texttt{.env} file, keeping sensitive information like API keys and database credentials out of source code.

\section{AI and Machine Learning Integration}

\subsection{Google Generative AI (Gemini API)}
The Gemini 1.5 Flash model serves as the core intelligence for fabric analysis. It is a multimodal large language model capable of understanding images and generating structured text responses. Key advantages:
\begin{itemize}
    \item Zero‑shot classification – no need for custom training datasets
    \item Ability to interpret fabric texture, color, and soil level from images
    \item Fast response times suitable for real‑time user interaction
\end{itemize}

\subsection{Pillow (PIL)}
Pillow is a Python imaging library used to open, manipulate, and convert image files before sending them to the Gemini API. It ensures that images are in a compatible format (JPEG/PNG) and can be resized if necessary to optimize API call performance.

\subsection{Tenacity}
Tenacity is a retry library used to implement automatic retries with exponential backoff when the Gemini API returns rate‑limit errors (\texttt{ResourceExhausted}). This improves system robustness under high load.

\section{Mobile Application Technologies}

\subsection{React Native with Expo}
React Native is a framework for building native mobile applications using JavaScript and React. Expo is a set of tools and services built around React Native that simplifies development, testing, and deployment.

Reasons for choosing React Native + Expo:
\begin{itemize}
    \item Cross‑platform (iOS and Android) from a single codebase
    \item Fast refresh for rapid UI iteration
    \item Access to device camera via Expo Camera module
    \item Large ecosystem of pre‑built components
\end{itemize}

\subsection{TypeScript}
TypeScript adds static typing to JavaScript, catching errors during development and improving code maintainability. Custom interfaces were defined for fabric data, wash parameters, and API responses.

\subsection{React Navigation}
React Navigation is used to implement tab‑based navigation (Dashboard, Scan, Progress, Override, Settings) with smooth transitions and state persistence.

\section{Hardware Components}

\subsection{Microcontroller: Arduino Uno WiFi R4}
The Arduino Uno WiFi R4 is the brain of the hardware system. It features:
\begin{itemize}
    \item Built‑in WiFi module for communication with the backend APIs
    \item Sufficient digital I/O pins to control relays, pumps, and motor drivers
    \item 5V logic level compatible with most modules
    \item Arduino ecosystem with extensive libraries
\end{itemize}

\subsection{Motor Control Relays}
Two 240V AC, 10A rated relay modules are used to switch the washing machine’s three-phase motor. One relay controls clockwise rotation, the other counter-clockwise. A third relay (5A) controls the water inlet solenoid.

\subsection{Detergent Dispensing System}
\begin{itemize}
    \item \textbf{5V DC Water Pump:} Draws liquid detergent from a container.
    \item \textbf{TAS122 MOSFET Transistor:} Acts as a high‑speed switch controlled by PWM from the Arduino.
    \item \textbf{1N4007 Flyback Diode:} Protects the transistor from voltage spikes when the pump is turned off.
\end{itemize}

\subsection{Water Inlet Solenoid Valve}
A 240V AC normally‑closed solenoid valve controls the flow of water from the household supply into the washing machine drum. It is activated by Relay 3.

\subsection{Water Outlet Mechanism}
\begin{itemize}
    \item \textbf{L298N Motor Driver:} Controls a 12V DC motor that opens and closes the drain valve.
    \item \textbf{12V DC Motor:} Mechanically linked to the drain valve.
    \item \textbf{12V Power Supply:} Dedicated supply for the motor driver.
\end{itemize}

\subsection{Power Supplies}
\begin{itemize}
    \item 5V supply for Arduino and relay logic (via USB or external adapter)
    \item 12V supply for L298N and outlet motor
    \item 240V AC mains for the washing machine motor and solenoid valve (switched via relays)
\end{itemize}

\section{Libraries and Dependencies}

\subsection{Python Backend}
\begin{table}[H]
    \centering
    \begin{tabular}{|l|l|}
        \hline
        \textbf{Library} & \textbf{Purpose} \\
        \hline
        fastapi & Web framework \\
        \hline
        uvicorn & ASGI server \\
        \hline
        pydantic & Data validation \\
        \hline
        google-generativeai & Gemini API client \\
        \hline
        Pillow & Image processing \\
        \hline
        python-dotenv & Environment variable management \\
        \hline
        tenacity & Retry logic \\
        \hline
        httpx & Async HTTP client (for internal API calls) \\
        \hline
    \end{tabular}
    \caption{Python Backend Dependencies}
    \label{tab:py_deps}
\end{table}

\subsection{Arduino (C++)}
\begin{table}[H]
    \centering
    \begin{tabular}{|l|l|}
        \hline
        \textbf{Library} & \textbf{Purpose} \\
        \hline
        WiFi.h & WiFi connectivity \\
        \hline
        ArduinoJson.h & Parsing JSON parameters from API \\
        \hline
    \end{tabular}
    \caption{Arduino Libraries}
    \label{tab:arduino_deps}
\end{table}

\subsection{React Native}
\begin{table}[H]
    \centering
    \begin{tabular}{|l|l|}
        \hline
        \textbf{Library} & \textbf{Purpose} \\
        \hline
        expo & Core Expo framework \\
        \hline
        expo-camera & Camera access for fabric scanning \\
        \hline
        expo-status-bar & Status bar management \\
        \hline
        @react-navigation/native & Navigation \\
        \hline
        @react-navigation/bottom-tabs & Tab navigation \\
        \hline
        axios & HTTP requests to backend APIs \\
        \hline
        react-native-svg & Icons and progress indicators \\
        \hline
    \end{tabular}
    \caption{React Native Dependencies}
    \label{tab:rn_deps}
\end{table}

\section{Testing Tools}
\begin{table}[H]
    \centering
    \begin{tabular}{|l|l|}
        \hline
        \textbf{Tool} & \textbf{Purpose} \\
        \hline
        Postman & Manual testing of API endpoints, debugging responses \\
        \hline
        pytest & Unit testing for Python backend functions (planned) \\
        \hline
        Arduino Serial Monitor & Debugging Arduino code, viewing real‑time logs \\
        \hline
        Expo Go & Live testing of mobile app on physical devices \\
        \hline
        Browser Developer Tools & Inspecting network requests, console logs \\
        \hline
    \end{tabular}
    \caption{Testing Tools}
    \label{tab:test_tools}
\end{table}

\section{Version Control and Collaboration}
\begin{table}[H]
    \centering
    \begin{tabular}{|l|l|}
        \hline
        \textbf{Tool} & \textbf{Purpose} \\
        \hline
        Git & Distributed version control \\
        \hline
        GitHub & Remote repository hosting, issue tracking, collaboration \\
        \hline
        .gitignore & Excluding sensitive files (\texttt{.env}, \texttt{node\_modules}) \\
        \hline
    \end{tabular}
    \caption{Version Control}
    \label{tab:vc}
\end{table}

\section{Deployment and Hosting}
\begin{table}[H]
    \centering
    \begin{tabular}{|l|l|}
        \hline
        \textbf{Component} & \textbf{Deployment Strategy} \\
        \hline
        FastAPI Backend & Can be deployed on any cloud VM (AWS EC2, DigitalOcean) or locally on a Raspberry Pi near the washing machine \\
        \hline
        Mobile App & Published via Expo Application Services (EAS) or standalone builds for App Store / Google Play \\
        \hline
        Arduino Firmware & Flashed via USB; updates can be delivered over‑the‑air using WiFi (future enhancement) \\
        \hline
    \end{tabular}
    \caption{Deployment and Hosting}
    \label{tab:deploy}
\end{table}

\section{Security Measures}
\begin{table}[H]
    \centering
    \begin{tabular}{|l|l|}
        \hline
        \textbf{Measure} & \textbf{Implementation} \\
        \hline
        API Key Protection & Stored in environment variables, never hard‑coded \\
        \hline
        Rate Limiting & Implemented via tenacity retries and future middleware \\
        \hline
        Input Validation & Pydantic models ensure only valid data is accepted \\
        \hline
        Network Security & WiFi secured with WPA2; local network communication assumed safe for prototype \\
        \hline
        Physical Safety & Relays fail open; software watchdogs; emergency stop in mobile app \\
        \hline
    \end{tabular}
    \caption{Security Measures}
    \label{tab:security}
\end{table}
